\chapter{راه‌اندازی ارتباط و اندازه‌گیری با \lr{ADE7880}}

\section{دامنه این مرحله و نتیجه قابل انتظار}
در این مرحله رابط
\lr{SPI2}
میکروکنترلر دوباره فعال، کتابخانه
\lr{ADE7880}
مستقل از خانواده میکروکنترلر بازنویسی و خواندن دوره‌ای کمیت‌های اصلی به برنامه افزوده شد. برنامه در هر نمونه، ولتاژ مؤثر، جریان مؤثر، توان اکتیو، توان ظاهری و ضریب توان هر سه فاز و همچنین جریان مؤثر نول را می‌خواند. داده خام کامل در متغیر
\lr{g\_app\_electrical}
منتشر می‌شود و در پنجره
\lr{Watch}
قابل مشاهده است. علاوه بر داده خام، مقدارهای مهندسی هر سه فاز با جدول ضرایب پیش‌فرض برد در آرایه‌های متناظر
\lr{voltage\_v}
قرار می‌گیرد. معتبر بودن هر خانه و منبع ضریب آن به‌ترتیب در
\lr{voltage\_valid}
و
\lr{voltage\_source}
ثبت می‌شود.

خواندن رجیستر با تبدیل مقدار خام به واحد فیزیکی یکسان نیست. ضریب توان از قالب ثابت خود تراشه مستقیماً به بازه منفی یک تا مثبت یک تبدیل‌پذیر است، اما تبدیل دیگر کمیت‌ها به ولت، آمپر و وات وابسته به نسبت ترانس جریان، شبکه تقسیم ولتاژ، مقاومت بار، بهره ورودی و کالیبراسیون واقعی همان برد است. جدول پیش‌فرض برای راه‌اندازی و مشاهده هر سه فاز مفید است، ولی تا آزمون با مرجع معتبر نباید مبنای حفاظت، ایمنی یا اندازه‌گیری نهایی قرار گیرد.

\section{منبع بازیابی اتصال‌های قبلی}
مخزن هنگام آغاز این مرحله هیچ تاریخچه
\lr{Git}
نداشت، ولی فایل پشتیبان قدیمی همین پروژه در مسیر بیرونی
\filepath{embedded/software/backup/Main/STM32F107.ioc}
هنوز موجود بود. تنظیمات آن با گزارش قبلی و سورس قدیمی تطبیق داده شد. هر سه منبع اتصال زیر را تأیید کردند:

\begin{longtable}{p{0.17\textwidth}p{0.25\textwidth}p{0.48\textwidth}}
\toprule
پایه میکرو & سیگنال & تنظیم لازم \\
\midrule
\endhead
\lr{PB12} & \lr{ADE\_CS / SS} & خروجی پوش‌پول، فعال‌پایین، سطح بی‌کار بالا \\
\lr{PB13} & \lr{SPI2\_SCK} & خروجی عملکرد جایگزین پوش‌پول \\
\lr{PB14} & \lr{SPI2\_MISO} & ورودی بدون کشش داخلی \\
\lr{PB15} & \lr{SPI2\_MOSI} & خروجی عملکرد جایگزین پوش‌پول \\
\bottomrule
\end{longtable}

در فایل‌های قدیمی هیچ پایه‌ای برای
\lr{RESET}
یا
\lr{IRQ1}
تراشه تعریف نشده بود. به همین دلیل نسخه حاضر از بازنشانی نرم‌افزاری و بررسی رجیستر
\lr{STATUS1.RSTDONE}
استفاده می‌کند. اگر در شماتیک واقعی این دو خط به میکرو متصل‌اند، شماره پایه‌ها باید در مرحله بعد از روی شماتیک قطعی شود؛ تا آن زمان نباید پایه‌ای برای آن‌ها حدس زده شود.

\section{کنترل الزامات دیتاشیت}
مرجع این پیاده‌سازی دیتاشیت رسمی بازبینی
\lr{C}
شرکت
\lr{Analog Devices}
است:
\href{https://www.analog.com/media/en/technical-documentation/data-sheets/ADE7880.pdf}{\lr{ADE7880 Rev. C}}

نکات زیر مستقیماً در طراحی رعایت شده‌اند:

\begin{itemize}
  \item تراشه پس از روشن‌شدن تقریباً ۴۰ میلی‌ثانیه برای آماده‌شدن نیاز دارد؛ کتابخانه ۵۰ میلی‌ثانیه حاشیه می‌دهد.
  \item رابط پیش‌فرض پس از \lr{reset} از نوع \lr{I2C} است. سه گذار پایه \lr{SS} از بالا به پایین رابط را به \lr{SPI} تغییر می‌دهد.
  \item نوشتن رجیستر هشت‌بیتی \lr{CONFIG2} پس از انتخاب \lr{SPI}، انتخاب رابط را قفل می‌کند.
  \item داده روی لبه بالارونده \lr{SCLK} نمونه‌برداری و پرارزش‌ترین بیت ابتدا منتقل می‌شود.
  \item بیشینه کلاک مجاز \lr{SPI} برابر ۲.۵ مگاهرتز است. نرخ انتخاب‌شده ۵۶۲.۵ کیلوبیت‌برثانیه حاشیه کافی دارد.
  \item پایان \lr{reset} نرم‌افزاری با دو شرط کنترل می‌شود: \lr{SWRST=0} و \lr{RSTDONE=1}.
  \item پردازشگر داخلی تنها پس از نوشتن و بازخوانی موفق مقدار یک در رجیستر \lr{RUN} شروع می‌شود.
\end{itemize}

\section{تنظیم از صفر در \lr{STM32CubeMX}}
فایل
\filepath{STM32F107.ioc}
مرجع تنظیم مولد است. برای بازسازی دستی همین پیکربندی مراحل زیر انجام شود:

\begin{enumerate}
  \item فایل پروژه را در نسخه سازگار \lr{STM32CubeMX} باز کنید و قطعه \lr{STM32F107VCTx} با بسته \lr{LQFP100} را تغییر ندهید.
  \item از بخش \lr{Connectivity}، رابط \lr{SPI2} را روی \lr{Full-Duplex Master} قرار دهید.
  \item کنترل کنید \lr{PB13} برابر \lr{SPI2\_SCK}، \lr{PB14} برابر \lr{SPI2\_MISO} و \lr{PB15} برابر \lr{SPI2\_MOSI} شده باشد.
  \item پایه \lr{PB12} را \lr{GPIO Output} با برچسب \lr{ADE\_CS} انتخاب کنید. سطح اولیه باید \lr{High}، نوع خروجی \lr{Push-Pull}، کشش \lr{No pull} و سرعت \lr{High} باشد.
  \item در تنظیم‌های \lr{SPI2} حالت ارباب، جهت دوخطه، طول هشت بیت، \lr{CPOL=High}، \lr{CPHA=2 Edge}، انتخاب تراشه نرم‌افزاری، \lr{MSB First} و تقسیم‌کننده ۶۴ را انتخاب کنید.
  \item مقدار محاسبه‌شده کلاک باید با کلاک گذرگاه اول ۳۶ مگاهرتز برابر ۵۶۲.۵ کیلوبیت‌برثانیه باشد.
  \item فعال‌کردن وقفه \lr{SPI2} برای روش \lr{polling} این مرحله لازم نیست. افزودن بی‌دلیل وقفه فقط پیچیدگی ایجاد می‌کند.
  \item در \lr{Project Manager > Code Generator} گزینه نگه‌داری کد کاربر فعال و حذف فایل‌های قبلی غیرفعال باقی بماند.
  \item کد را تولید و سپس تفاوت‌ها را با \lr{Git} بازبینی کنید. وجود \lr{MX\_SPI2\_Init} و تنظیم پایه‌ها نتیجه مورد انتظار مولد است.
\end{enumerate}

فراخوانی‌های برنامه داخل ناحیه‌های حفاظت‌شده فایل اصلی هستند:

\begin{LTR}
\begin{lstlisting}[language=C]
/* USER CODE BEGIN 2 */
App_Init(&hspi2);
/* USER CODE END 2 */
/* USER CODE BEGIN WHILE */
while (1) {
    App_Process();
/* USER CODE END WHILE */
}
\end{lstlisting}
\end{LTR}

در نتیجه، تولید مجدد \lr{CubeMX} ممکن است بدنه‌های مولد
\lr{MX\_SPI2\_Init}
و
\lr{HAL\_SPI\_MspInit}
را بازسازی کند، اما منطق کتابخانه و برنامه در پوشه‌های
\filepath{Libraries/ADE7880}
و
\filepath{Application}
پاک نمی‌شود و دو نقطه اتصال نیز داخل بلوک کاربر حفظ می‌شوند.

\section{معماری کتابخانه قابل استفاده مجدد}
کتابخانه فقط دو فایل دارد:

\begin{itemize}
  \item \filepath{Libraries/ADE7880/Inc/ade7880.h}: انواع، رجیسترها و رابط عمومی؛
  \item \filepath{Libraries/ADE7880/Src/ade7880.c}: پروتکل، \lr{reset}، خواندن و کالیبراسیون.
\end{itemize}

کتابخانه هدر
\lr{stm32f1xx\_hal.h}
را نیاز ندارد. وابستگی سخت‌افزار با ساختار
\lr{ADE7880\_Transport}
به آن تزریق می‌شود:

\begin{LTR}
\begin{lstlisting}[language=C]
typedef struct {
    void *context;
    ADE7880_Status (*write)(void *, const uint8_t *, size_t, uint32_t);
    ADE7880_Status (*read)(void *, uint8_t *, size_t, uint32_t);
    void (*select)(void *, bool active);
    void (*delay_ms)(void *, uint32_t);
} ADE7880_Transport;
\end{lstlisting}
\end{LTR}

برای انتقال کتابخانه به میکرو دیگر فقط همین چهار \lr{callback} پیاده می‌شود. برای مثال در \lr{STM32}، \lr{callback} نوشتن به
\lr{HAL\_SPI\_Transmit}
و \lr{callback} خواندن به
\lr{HAL\_SPI\_Receive}
وصل شده است. در پروژه دیگر می‌توان آن‌ها را به درایور
\lr{LL}،
\lr{ESP-IDF}،
\lr{Zephyr}
یا یک \lr{SPI} نرم‌افزاری متصل کرد، بدون آن‌که فایل کتابخانه تغییر کند.

راهنمای مستقل و کوتاه استفاده از کتابخانه نیز در فایل
\filepath{Libraries/ADE7880/README.md}
قرار دارد؛ بنابراین برای انتقال ماژول به پروژه دیگر لازم نیست گزارش کامل محصول همراه آن کپی شود.

\section{قرارداد کامنت‌گذاری و روش خواندن کد}
تمام تعریف‌های عمومی کتابخانه در
\filepath{ade7880.h}
با قالب
\lr{Doxygen}
انگلیسی مستند شده‌اند. هر تابع عمومی باید هدف، پارامترهای ورودی، خروجی و خطاهای قابل انتظار را داشته باشد. اعضای ساختار، مقدارهای
\lr{enum}
و رجیسترها نیز کنار تعریف خود توضیح انگلیسی دارند. فایل پیاده‌سازی علاوه بر توضیح هر تابع، دلیل مراحل غیر بدیهی مانند ترتیب بایت، بسط علامت، جلوگیری از انتشار نمونه ناقص و پاک‌کردن
\lr{RSTDONE}
را ثبت می‌کند.

قاعده ادامه پروژه چنین است:

\begin{itemize}
  \item نام تابع و متغیر باید هدف را بیان کند و کامنت جای نام‌گذاری نامناسب را نگیرد؛
  \item رابط عمومی با \lr{@brief}، \lr{@param} و \lr{@return} انگلیسی مستند شود؛
  \item عدد ثابت سخت‌افزاری باید واحد و علت داشته باشد؛
  \item کامنت باید دلیل یا قرارداد را توضیح دهد، نه آن‌که همان دستور ساده را تکرار کند؛
  \item هر تغییر رفتار در گزارش فارسی و فایل وضعیت نیز ثبت شود.
\end{itemize}

\section{مرجع سریع توابع کتابخانه}
\begin{itemize}
  \item \lr{ADE7880\_Init}: ثبت و کنترل چهار \lr{callback} سخت‌افزار؛ هنوز با تراشه تبادل انجام نمی‌دهد.
  \item \lr{ADE7880\_Begin}: انتخاب \lr{SPI}، بازنشانی نرم‌افزاری، انتظار محدود و شروع \lr{DSP}.
  \item \lr{ADE7880\_ReadRegister}: خواندن عمومی رجیسترهای ۸، ۱۶ یا ۳۲بیتی.
  \item \lr{ADE7880\_WriteRegisterVerified}: نوشتن رجیستر و مقایسه بازخوانی برای تشخیص خطا.
  \item \lr{ADE7880\_ReadPhase}: خواندن همه کمیت‌های خام یک فاز.
  \item \lr{ADE7880\_ReadAll}: ساخت نمونه کامل سه فاز و نول؛ در خطا خروجی نیمه‌کاره منتشر نمی‌شود.
  \item \lr{ADE7880\_CalculateLinearCalibration}: محاسبه \lr{scale/offset} از دو نقطه خام و مرجع.
  \item \lr{ADE7880\_ApplyLinearCalibration}: تبدیل یک شمار خام به واحد مهندسی.
  \item \lr{ADE7880\_Convert}: تبدیل یک نمونه کامل با مجموعه ضرایب همه کانال‌ها.
  \item \lr{App\_SetPhaseCalibration}: جایگزینی مستقل ضریب ولتاژ، جریان، توان اکتیو یا توان ظاهری هر فاز.
  \item \lr{App\_CalibratePhaseMultiPoint}: برازش و نصب کالیبراسیون چندنقطه‌ای برای یک کمیت و یک فاز.
  \item \lr{App\_CalibrateNeutralCurrentMultiPoint}: برازش و نصب مستقل جریان نول.
  \item \lr{\small App\_ResetCalibrationToDefaults}: حذف ضرایب زمان اجرا و بازگرداندن جدول آغازکار.
\end{itemize}

ساده‌ترین مسیر استفاده در همین پروژه فقط دو فراخوانی دارد: یک بار
\lr{App\_Init(\&hspi2)}
بعد از راه‌اندازی \lr{SPI2} و سپس اجرای پیوسته
\lr{App\_Process()}
در حلقه اصلی. لایه برنامه ساختار انتقال، بازیابی خطا، زمان نمونه‌برداری و تبدیل ولتاژ را مدیریت می‌کند؛ کاربر معمولاً فقط
\lr{g\_app\_electrical}
را می‌خواند.

\section{توالی راه‌اندازی برنامه}
تابع
\lr{App\_Init}
پس از راه‌اندازی پایه‌ها و
\lr{SPI2}
فراخوانی می‌شود. توالی دقیق آن چنین است:

\begin{enumerate}
  \item ماژول دو \lr{LED} راه‌اندازی می‌شود.
  \item \lr{callback}های \lr{HAL} در یک \lr{ADE7880\_Transport} قرار می‌گیرند.
  \item \lr{ADE7880\_Init} اعتبار \lr{callback}ها را بررسی و \lr{CS} را غیرفعال می‌کند.
  \item \lr{ADE7880\_Begin} پنجاه میلی‌ثانیه صبر، سه پالس \lr{CS} ایجاد و \lr{SPI} را با \lr{CONFIG2} قفل می‌کند.
  \item \lr{reset} نرم‌افزاری با مهلت ۱۰۰ میلی‌ثانیه اجرا می‌شود؛ حلقه نامحدود وجود ندارد.
  \item رجیستر \lr{RUN} نوشته و بازخوانی می‌شود.
  \item رجیستر هشت‌بیتی \lr{VERSION} خوانده و در وضعیت دیباگ ثبت می‌شود.
  \item موفقیت حالت رقص \lr{LED} و خطا تناوب سریع دو \lr{LED} را فعال می‌کند.
\end{enumerate}

اگر ارتباط قطع باشد، برنامه در راه‌اندازی گیر نمی‌کند. وضعیت آفلاین ثبت و هر دو ثانیه تلاش مجدد انجام می‌شود. پس از بازیابی ارتباط، الگوی \lr{LED} خودکار به رقص برمی‌گردد. در حالت سالم یک \lr{snapshot} کامل تقریباً هر یک ثانیه خوانده می‌شود.

\section{رجیسترها و قالب داده خوانده‌شده}
\begin{longtable}{p{0.14\textwidth}p{0.16\textwidth}p{0.20\textwidth}p{0.36\textwidth}}
\toprule
کمیت & فاز \lr{A} & فازهای \lr{B/C} & تفسیر \\
\midrule
\endhead
جریان مؤثر & \lr{0x43C0} & \lr{0x43C2/0x43C4} & مقدار ۲۴بیتی بدون علامت کاربردی \\
ولتاژ مؤثر & \lr{0x43C1} & \lr{0x43C3/0x43C5} & مقدار ۲۴بیتی بدون علامت کاربردی \\
توان اکتیو & \lr{0xE513} & \lr{0xE514/0xE515} & مقدار ۲۴بیتی علامت‌دار با بسط علامت \\
توان ظاهری & \lr{0xE519} & \lr{0xE51A/0xE51B} & داده ۲۴بیتی در کلمه ۳۲بیتی \\
ضریب توان & \lr{0xE902} & \lr{0xE903/0xE904} & عدد علامت‌دار \lr{Q1.15}؛ هر \lr{LSB} برابر \lr{$2^{-15}$} \\
جریان نول & \lr{0x43C6} & ندارد & مقدار ۲۴بیتی \\
\bottomrule
\end{longtable}

در نسخه قدیمی ضریب توان از رجیستر هارمونیک مشترک
\lr{FPF}
خوانده می‌شد و انتخاب فاز نیز درست تضمین نشده بود. نسخه حاضر مستقیماً از سه رجیستر
\lr{APF/BPF/CPF}
استفاده می‌کند؛ بنابراین هم ساده‌تر و هم مطابق بخش \lr{Power Factor} دیتاشیت است.

\section{خواندن در \lr{VS Code} و \lr{PlatformIO}}
پس از بازکردن ریشه پروژه، افزونه
\lr{PlatformIO IDE}
را نصب و محیط
\lr{stm32f107vc\_jlink}
را انتخاب کنید. ساخت از ترمینال:

\begin{LTR}
\begin{verbatim}
pio run -e stm32f107vc_jlink
\end{verbatim}
\end{LTR}

اگر فرمان
\lr{pio}
در \lr{PATH} نیست، روی ویندوز می‌توان فایل اجرایی محیط خود \lr{PlatformIO} را اجرا کرد:

\begin{LTR}
\begin{verbatim}
%USERPROFILE%\.platformio\penv\Scripts\platformio.exe run \
  -e stm32f107vc_jlink
\end{verbatim}
\end{LTR}

برای پروگرام:

\begin{LTR}
\begin{verbatim}
pio run -e stm32f107vc_jlink --target upload
\end{verbatim}
\end{LTR}

برای دیباگ، \lr{breakpoint} را بعد از انتساب
\lr{g\_app\_electrical.raw}
در
\lr{App\_Process}
قرار دهید و متغیر
\lr{g\_app\_electrical}
را به \lr{Watch} اضافه کنید. نتیجه سالم باید چنین ویژگی‌هایی داشته باشد:

\begin{itemize}
  \item \lr{online=true}؛
  \item \lr{last\_status=ADE7880\_STATUS\_OK}؛
  \item \lr{die\_version} دارای مقدار ثابت و تکرارپذیر؛
  \item افزایش \lr{successful\_samples} تقریباً هر ثانیه؛
  \item تازه‌شدن \lr{updated\_at\_ms}؛
  \item ثابت‌ماندن \lr{communication\_errors} در صفر.
\end{itemize}

\section{ساخت و دیباگ در \lr{Keil MDK}}
فایل
\filepath{MDK-ARM/STM32F107.uvprojx}
باز شود. کتابخانه از قبل در گروه
\lr{Libraries/Reusable}
عضو است و مسیر \lr{include} آن نیز ثبت شده است. بعد از تولید مجدد \lr{CubeMX} موارد زیر کنترل شود:

\begin{enumerate}
  \item فایل‌های \filepath{app.c} و \filepath{diagnostic_led.c} از پوشه \filepath{Application/Src} هنوز عضو پروژه باشند.
  \item فایل \filepath{Libraries/ADE7880/Src/ade7880.c} فقط یک بار عضو پروژه باشد.
  \item مسیرهای \filepath{Application/Inc} و \filepath{Libraries/ADE7880/Inc} در \lr{Include Paths} باشند.
  \item نمادهای \lr{USE\_HAL\_DRIVER} و \lr{STM32F107xC} تعریف شده باشند.
  \item هدف با گزینه \lr{Build Target} بدون \lr{warning} ساخته شود.
\end{enumerate}

در پنجره \lr{Watch} همان متغیر
\lr{g\_app\_electrical}
اضافه شود. برای جلوگیری از توقف طولانی سیستم قدرت، \lr{breakpoint} تنها هنگام جدا بودن بارها استفاده شود. اگر برنامه باید آزاد اجرا شود، از \lr{Watch} با به‌روزرسانی دوره‌ای و بدون توقف مکرر بهره بگیرید.

ساخت واقعی همین پروژه با \lr{Keil} و \lr{Arm Compiler 6.24} اجرا شد و با نتیجه
\lr{0 Error(s), 0 Warning(s)}
پایان یافت. اندازه‌های گزارش‌شده برابر
\lr{Code=9628}،
\lr{RO-data=476}،
\lr{RW-data=12}
و
\lr{ZI-data=1460}
بایت بودند و فایل \lr{HEX} نیز تولید شد.

خطاهای اولیه همگی از ناشناخته‌بودن
\lr{DMA\_HandleTypeDef}
آغاز می‌شدند. فایل پیاده‌سازی \lr{DMA} مربوط به \lr{HAL} در پروژه وجود داشت، ولی تعریف
\lr{HAL\_DMA\_MODULE\_ENABLED}
در
\filepath{Core/Inc/stm32f1xx_hal_conf.h}
غیرفعال بود. با فعال‌کردن این تعریف، نوع‌های مورد نیاز \lr{handle}های
\lr{I2C/SPI/TIM/UART}
پیش از استفاده در هدرهای \lr{HAL} تعریف شدند. این یک وابستگی تعریفی در \lr{HAL} است و به‌تنهایی به معنی استفاده منطق برنامه از انتقال \lr{DMA} نیست.

\section{کالیبراسیون دو نقطه‌ای}
نوع
\lr{ADE7880\_LinearCalibration}
دو عدد \lr{scale} و \lr{offset} دارد و رابطه زیر را اعمال می‌کند:

\begin{center}
\lr{engineering = raw * scale + offset}
\end{center}

برای نمونه، فرض کنید در دو نقطه مرجع ۱۰۰ و ۲۳۰ ولت، رجیستر ولتاژ به‌ترتیب ۱۱۰۰۰۰۰ و ۲۵۵۰۰۰۰ باشد. محاسبه ضریب بدون دست‌کاری دستی فرمول چنین است:

\Needspace{13\baselineskip}
\begin{LTR}
\begin{lstlisting}[language=C]
ADE7880_LinearCalibration voltage_a;

if (ADE7880_CalculateLinearCalibration(
        1100000.0F, 100.0F,
        2550000.0F, 230.0F,
        &voltage_a) == ADE7880_STATUS_OK) {
    App_SetVoltageCalibration(ADE7880_PHASE_A, &voltage_a);
}
\end{lstlisting}
\end{LTR}

سپس
\lr{ADE7880\_Convert}
یک \lr{snapshot} خام و ساختار ضرایب را به مقادیر مهندسی تبدیل می‌کند. بهتر است ضرایب هر فاز مستقل باشند. روند عملی کالیبراسیون:

\begin{enumerate}
  \item برد و مرجع اندازه‌گیری حداقل تا رسیدن دما و تغذیه به حالت پایدار روشن بمانند.
  \item برای هر فاز دو یا چند نقطه دور از هم انتخاب شود؛ نقطه صفر به‌تنهایی برای \lr{RMS} مناسب نیست.
  \item در هر نقطه میانگین چند ده \lr{snapshot} خام ثبت شود.
  \item مقدار مرجع کالیبره‌شده هم‌زمان ثبت و \lr{scale/offset} محاسبه شود.
  \item ضرایب با شماره سریال برد، تاریخ، ابزار مرجع و شرایط آزمون ذخیره شوند.
  \item پس از اعمال ضرایب، نقطه سوم که در محاسبه استفاده نشده کنترل شود.
\end{enumerate}

\subsection{ضرایب ولتاژ بازیابی‌شده و خروجی فعلی}
در سورس قدیمی همین برد تبدیل زیر پیدا شد:

\begin{center}
\lr{phase A volts = AVRMS * 0.00055963 * 1.017}

\lr{phase B volts = BVRMS * 0.00055963}
\end{center}

این دو ضریب اکنون هنگام
\lr{App\_Init}
بارگذاری می‌شوند. پس از نخستین نمونه موفق، خروجی‌های زیر در
\lr{Watch}
قابل استفاده‌اند:

\begin{itemize}
  \item \lr{voltage\_v[0/1/2]}: ولتاژ تقریبی فازهای \lr{A/B/C}؛
  \item \lr{voltage\_valid[0/1/2]=true}: تبدیل جدول پیش‌فرض پس از نخستین نمونه کامل انجام شده است؛
  \item \lr{\scriptsize voltage\_source[0/1/2]}: مقدار \lr{\scriptsize APP\_CALIBRATION\_DEFAULT} تا زمانی که ضریب اندازه‌گیری‌شده نصب نشده باشد؛
  \item ضریب اولیه فاز \lr{C} از فاز \lr{B} کپی شده و باید روی سخت‌افزار واقعی کنترل شود.
\end{itemize}

با فراخوانی
\lr{App\_SetVoltageCalibration}
منبع ضریب همان فاز «کاربر» ثبت می‌شود و نمونه بعدی با ضریب جدید منتشر خواهد شد. برای تکمیل کالیبراسیون جریان و توان هنوز نسبت \lr{CT}، مقدار مقاومت \lr{burden}، شبکه ورودی و محدوده نامی لازم است.

\subsection{مشاهده جریان، توان و ضریب توان در دیباگر}
درایور از ابتدا رجیسترهای جریان \lr{RMS}، توان اکتیو، توان ظاهری و ضریب توان را نیز در هر \lr{snapshot} می‌خواند. این داده‌ها در ساختار خام زیر قرار دارند:

\begin{itemize}
  \item \lr{raw.phase[0]} برای فاز \lr{A}؛
  \item \lr{raw.phase[1]} برای فاز \lr{B}؛
  \item \lr{raw.phase[2]} برای فاز \lr{C}؛
  \item \lr{raw.neutral\_current\_rms} برای جریان نول.
\end{itemize}

برای نمونه، پیش از انجام هر کالیبراسیون می‌توان در \lr{Watch} ساختار
\lr{g\_app\_electrical.raw.phase[0]}
را باز کرد و اعضای
\lr{current\_rms}،
\lr{active\_power}،
\lr{apparent\_power}
و
\lr{power\_factor\_q15}
را برای فاز \lr{A} دید. این مقادیر به‌جز ضریب توان، شمار خام پردازشگر داخلی تراشه هستند و واحد آمپر یا وات ندارند.

نسخه جدید برنامه فیلدهای مستقیم و خوانای زیر را نیز منتشر می‌کند:

\begin{longtable}{p{0.38\textwidth}p{0.52\textwidth}}
\toprule
فیلد در \lr{g\_app\_electrical} & معنی و شرط استفاده \\
\midrule
\endhead
\lr{current\_a} & آرایه جریان \lr{RMS} برحسب آمپر. \\
\lr{active\_power\_w} & آرایه توان اکتیو برحسب وات. \\
\lr{apparent\_power\_va} & آرایه توان ظاهری برحسب ولت‌آمپر. \\
\lr{power\_factor} & آرایه ضریب توان علامت‌دار در بازه ۱- تا ۱+. \\
\lr{neutral\_current\_a} & جریان نول برحسب آمپر؛ فقط با پرچم اعتبار متناظر قابل استفاده است. \\
\bottomrule
\end{longtable}

آرایه \lr{power\_factor\_abs} بزرگی ضریب توان را در بازه صفر تا یک و آرایه \lr{current\_polarity} جهت جبران \lr{CT} هر فاز را به‌صورت \lr{NORMAL} یا \lr{REVERSED} نشان می‌دهد.

سه عضو هر آرایه به‌ترتیب فازهای \lr{A/B/C} هستند. برای جریان، توان اکتیو و توان ظاهری به‌ترتیب پرچم‌های زیر کنترل شوند:

\begin{itemize}
  \item \lr{current\_valid}؛
  \item \lr{active\_power\_valid}؛
  \item \lr{apparent\_power\_valid}.
\end{itemize}

ضریب توان پس از یک \lr{snapshot} کامل با \lr{power\_factor\_valid=true} قابل استفاده است.

رجیسترهای
\lr{APF/BPF/CPF}
اعداد علامت‌دار \lr{Q1.15} هستند؛ بنابراین هر \lr{LSB} وزن
\lr{$2^{-15}$}
دارد و برنامه آن‌ها را با تقسیم بر ۳۲۷۶۸ به ضریب توان تبدیل می‌کند. این تبدیل به نسبت \lr{CT} یا شبکه ولتاژ وابسته نیست. علامت \lr{PF} مطابق قرارداد تراشه اطلاعات پیش‌فاز/پس‌فاز بودن را نیز حمل می‌کند.

مقدار آمپر و وات به سخت‌افزار همان برد وابسته است. در نسخه نخست این بخش، چون ضریب کالیبراسیون جریان و توان بارگذاری نشده بود، برنامه این فیلدها را عمداً صفر و پرچم اعتبارشان را نادرست نگه می‌داشت؛ بنابراین تغییر داده خام همراه با صفر ماندن مقدار مهندسی نشانه خرابی \lr{SPI} نبود.

در بازبینی سورس پشتیبان همین برد، ضرایب زیر برای فازهای \lr{A/B} پیدا شدند. کد قدیمی نتیجه را در متغیر صحیح ۱۶بیتی و بسته ارتباطی قرار می‌داد. آزمون واقعی با لامپ حدود ۲۰ وات نشان داد خروجی قدیمی جریان برحسب صدم آمپر و خروجی توان برحسب دهم وات بوده است: عددهای \lr{17} و \lr{231} به‌ترتیب حدود \lr{0.17 A} و \lr{23.1 W} هستند. بنابراین تبدیل درست به واحدهای \lr{SI} چنین است:

\begin{center}
\lr{current amperes = IRMS * (0.00036565 / 100)}

\lr{active watts = WATT * (0.0170276 / 10)}
\end{center}

فاز \lr{A} در توان اصلاح اضافی \lr{1.02} داشت. این ضرایب اکنون هنگام
\lr{App\_Init}
بارگذاری می‌شوند. برای کامل‌بودن راه‌اندازی سه‌فاز، فاز \lr{C} فعلاً از مقادیر فاز \lr{B} استفاده می‌کند و جریان نول از ضریب جریان فازها آغاز می‌شود؛ بنابراین پس از نخستین نمونه موفق همه کانال‌ها خروجی مهندسی و پرچم اعتبار دارند. همه ضرایب آغازکار با
\lr{APP\_CALIBRATION\_DEFAULT}
علامت‌گذاری می‌شوند. این پرچم فقط موفقیت تبدیل را نشان می‌دهد و تأیید دقت ضریب جایگزین موقت نیست.

کد قدیمی توان ظاهری را به واحد مهندسی تبدیل نمی‌کرد. دیتاشیت \lr{ADE7880} و راهنمای \lr{AN-1171} بیان می‌کنند بهره توان‌های اکتیو و ظاهری هر فاز در تراشه با هم تطبیق داده شده است؛ ازاین‌رو مقدار اولیه \lr{VA} هر فاز با ضریب توان همان فاز آغاز می‌شود. وجود خانه مستقل \lr{apparent\_power[phase]} اجازه می‌دهد بعداً بدون اثر روی توان اکتیو، همان کانال جداگانه کالیبره شود.

برای تشخیص منشأ هر خروجی، آرایه‌های
\lr{voltage\_source}،
\lr{current\_source}،
\lr{active\_power\_source}
و
\lr{apparent\_power\_source}
در \lr{Watch} قابل مشاهده‌اند. مقدار
\lr{APP\_CALIBRATION\_USER}
یعنی ضریب پیش‌فرض همان خانه با کالیبراسیون اندازه‌گیری‌شده جایگزین شده است.

\subsection{جهت \lr{CT}، علامت توان و علامت ضریب توان}
اگر جهت اولیه یا ثانویه \lr{CT} برعکس شود، بردار جریان ۱۸۰ درجه تغییر می‌کند. در نتیجه علامت توان اکتیو و \lr{PF} علامت‌دار عوض می‌شود، ولی جریان \lr{RMS}، توان ظاهری و بزرگی \lr{PF} نباید تغییر کنند. برای نصب فقط مصرف‌کننده، جهت درست جهتی است که در هنگام روشن‌بودن بار، توان اکتیو مثبت باشد. اگر روی \lr{CT} علامت وجود دارد، معمولاً \lr{P1/K} سمت منبع، \lr{P2/L} سمت بار، \lr{S1/k} سمت ورودی جریان مثبت و \lr{S2/l} سمت ورودی منفی قرار می‌گیرد؛ نام نهایی پایه‌ها حتماً باید با شماتیک همان برد و دیتاشیت سازنده \lr{CT} تطبیق داده شود.

علامت \lr{PF} را نباید برای انتخاب جهت \lr{CT} به‌تنهایی استفاده کرد. طبق قرارداد \lr{ADE7880}، علامت این رجیستر اطلاعات پیش‌فاز/پس‌فاز را نیز حمل می‌کند. بنابراین یک لامپ \lr{LED} با ورودی یکسوساز و خازن می‌تواند با وجود توان مصرفی مثبت، \lr{PF} علامت‌دار منفی داشته باشد. فیلد
\lr{power\_factor}
علامت را حفظ می‌کند و فیلد جدید
\lr{power\_factor\_abs}
فقط بزرگی صفر تا یک را برای نمایش معمول ارائه می‌دهد.

اگر اصلاح فیزیکی جهت ممکن نباشد، پس از \lr{App\_Init} یک‌بار تنظیم زیر انجام می‌شود:

\begin{LTR}
\begin{lstlisting}[language=C]
(void)App_SetCurrentPolarity(
    ADE7880_PHASE_A,
    APP_CURRENT_POLARITY_REVERSED);
\end{lstlisting}
\end{LTR}

این تنظیم علامت توان اکتیو و \lr{PF} علامت‌دار را اصلاح می‌کند؛ جریان \lr{RMS}، توان ظاهری و قدرمطلق ضریب توان بدون تغییر می‌مانند. تشخیص و وارونه‌سازی خودکار در هر نمونه انجام نمی‌شود، زیرا در یک سامانه دوطرفه توان منفی می‌تواند صادرات واقعی باشد. مقدار انتخاب‌شده در آرایه
\lr{current\_polarity}
قابل مشاهده است و باید یک‌بار هنگام راه‌اندازی هر فاز تعیین شود.

\textbf{هشدار ایمنی:} ثانویه \lr{CT} در حالی که جریان از اولیه عبور می‌کند هرگز باز نشود؛ امکان ایجاد ولتاژ خطرناک وجود دارد. تغییر سیم‌بندی فقط پس از بی‌برق‌کردن اولیه یا با روش اتصال‌کوتاه ایمن مورد تأیید سازنده انجام شود.

\subsection{واحد مهندسی و نمایش سازگار با برنامه قدیمی}
برای جلوگیری از ابهام، فیلدهای محاسباتی همیشه واحد مهندسی صریح دارند:

\begin{itemize}
  \item \lr{current\_a}: آمپر؛
  \item \lr{active\_power\_w}: وات؛
  \item \lr{apparent\_power\_va}: ولت‌آمپر.
\end{itemize}

سه فیلد مستقل زیر فقط برای نمایش و پروتکل‌های \lr{fixed-point} در نظر گرفته شده‌اند:

\begin{itemize}
  \item \lr{current\_display}: پیش‌فرض صدم آمپر؛ بنابراین \lr{0.17 A} به‌شکل \lr{17} دیده می‌شود؛
  \item \lr{active\_power\_display}: پیش‌فرض دهم وات؛ بنابراین \lr{23.1 W} به‌شکل \lr{231} دیده می‌شود؛
  \item \lr{apparent\_power\_display}: پیش‌فرض دهم ولت‌آمپر.
\end{itemize}

دو ثابت زیر در \lr{app.h} پیش‌فرض را مشخص می‌کنند؛ مقدارهای آن‌ها به‌ترتیب ۱۰۰ و ۱۰ هستند:

\begin{LTR}
\begin{lstlisting}
APP_DEFAULT_CURRENT_DISPLAY_UNITS_PER_AMP
APP_DEFAULT_POWER_DISPLAY_UNITS_PER_WATT
\end{lstlisting}
\end{LTR}

تغییر موقت در زمان اجرا نیز بدون دست‌زدن به کالیبراسیون ممکن است:

\begin{LTR}
\begin{lstlisting}[language=C]
/* Make optional display fields use direct A and W units. */
(void)App_SetDisplayUnits(1.0F, 1.0F);
\end{lstlisting}
\end{LTR}

\subsection{جدول واحد ضرایب پیش‌فرض هر سه فاز}
همه ضرایب آغازکار برد در یک متغیر ثابت و خوانا به نام
\lr{default\_calibration}
در ابتدای فایل \lr{app.c}
قرار گرفته‌اند. ساختار هر خانه دقیقاً
\lr{\{scale, offset\}}
و رابطه تبدیل برای همه کمیت‌ها یکسان است:

\begin{center}
\lr{engineering = raw * scale + offset}
\end{center}

\begin{longtable}{p{0.20\textwidth}p{0.14\textwidth}p{0.31\textwidth}p{0.20\textwidth}}
\toprule
کمیت & کانال & \lr{scale} پیش‌فرض & \lr{offset} \\
\midrule
\endhead
ولتاژ \lr{RMS} & \lr{A} & \lr{0.00055963 * 1.017} & صفر \\
ولتاژ \lr{RMS} & \lr{B} & \lr{0.00055963} & صفر \\
ولتاژ \lr{RMS} & \lr{C} & \lr{0.00055963}، موقتاً کپی \lr{B} & صفر \\
جریان \lr{RMS} & \lr{A/B/C} & \lr{0.00036565 / 100} & صفر \\
توان اکتیو & \lr{A} & \lr{(0.0170276 / 10) * 1.02} & صفر \\
توان اکتیو & \lr{B/C} & \lr{0.0170276 / 10} & صفر \\
توان ظاهری & \lr{A/B/C} & برابر ضریب اولیه توان اکتیو همان فاز & صفر \\
جریان نول & \lr{N} & \lr{0.00036565 / 100}، موقتاً کپی جریان فاز & صفر \\
\bottomrule
\end{longtable}

این جدول عمداً تک‌تک کانال‌ها را نشان می‌دهد تا تغییر یک فاز به تغییر ناخواسته فاز دیگر منجر نشود. با وجود برابر بودن بعضی مقدارهای اولیه، هر خانه حافظه مستقل خود را دارد. تابع
\lr{App\_Init}
جدول را با
\lr{App\_ResetCalibrationToDefaults}
بارگذاری می‌کند. هر فراخوانی
\lr{App\_SetPhaseCalibration}
یا
\lr{App\_CalibratePhaseMultiPoint}
فقط خانه انتخاب‌شده را با منبع
\lr{APP\_CALIBRATION\_USER}
جایگزین می‌کند. فراخوانی مجدد بازنشانی نرم‌افزاری کالیبراسیون، تمام ضرایب زمان اجرای حافظه را حذف می‌کند:

\begin{LTR}
\begin{lstlisting}[language=C]
/* Discard runtime coefficients and restore the full startup table. */
App_ResetCalibrationToDefaults();
\end{lstlisting}
\end{LTR}

وجود ضریب جایگزین موقت برای فاز استفاده‌نشده هیچ ورودی را فعال یا خطر الکتریکی ایجاد نمی‌کند؛ فقط اگر بعداً سیگنال به آن فاز وصل شد، برنامه بلافاصله مقدار تقریبی منتشر می‌کند. پیش از اتکا، همان فاز مستقل با مرجع کالیبره می‌شود.

\subsection{استفاده تک‌فاز، دو کانال مستقل و سه‌فاز}
اندیس و رابط برنامه‌نویسی در هر سه حالت ثابت‌اند. محصول تک‌فاز می‌تواند فقط خانه
\lr{\small ADE7880\_PHASE\_A}
را مصرف کند. دو مصرف‌کننده مستقل می‌توانند از \lr{A} و \lr{B} استفاده کنند و کالیبراسیون هرکدام جدا باشد. دستگاه سه‌فاز هر سه خانه \lr{A/B/C} را پردازش می‌کند. نمونه تنظیم مستقل دو کانال:

\begin{LTR}
\begin{lstlisting}[language=C]
ADE7880_LinearCalibration current_a = {SCALE_A, OFFSET_A};
ADE7880_LinearCalibration current_b = {SCALE_B, OFFSET_B};

(void)App_SetPhaseCalibration(APP_PHASE_QUANTITY_CURRENT_RMS,
                              ADE7880_PHASE_A, &current_a);
(void)App_SetPhaseCalibration(APP_PHASE_QUANTITY_CURRENT_RMS,
                              ADE7880_PHASE_B, &current_b);
/* Phase C keeps its startup default until explicitly replaced. */
\end{lstlisting}
\end{LTR}

در سطح کتابخانه قابل‌حمل نیز
\lr{\small ADE7880\_ReadPhase}
برای خواندن انتخابی یک فاز و
\lr{\small ADE7880\_ReadAll}
برای نمونه اتمیک \lr{A/B/C} و نول موجود است. درایور هیچ نمونه سراسری از تراشه ندارد؛ بنابراین پروژه دیگر می‌تواند برای بیش از یک \lr{ADE7880}، یک
\lr{ADE7880\_Device}
مستقل با رابط انتقال مستقل بسازد.

\subsection{تصمیم کالیبراسیون: یک خط چندنقطه‌ای، نه ضرایب قطعه‌ای}
مسیر اندازه‌گیری \lr{ADE7880} ذاتاً خطی طراحی شده است. قرار دادن یک ضریب برای بار کم و ضریب دیگر برای بار زیاد می‌تواند در مرز دو بازه پرش بسازد و عیب واقعی \lr{offset}، نویز یا خطای فاز را پنهان کند. روش نخست این پروژه یک رابطه
\lr{engineering = raw * scale + offset}
است که با کمترین‌مربعات روی همه نقاط ثبت‌شده برازش می‌شود. حداقل دو نقطه لازم است، سه نقطه برای کم/میانی/زیاد حداقل عملی و پنج نقطه برای ارزیابی بهتر خطی‌بودن توصیه می‌شود.

راهنمای رسمی \lr{AN-1171} حداکثر سه مرحله متفاوت را معرفی می‌کند: \lr{gain}، \lr{phase} و \lr{offset}. این عبارت به معنی سه ضریب قطعه‌ای نیست. \lr{gain} برای هر دستگاه لازم است؛ اصلاح فاز برای \lr{CT}، به‌خصوص در ضریب توان پایین، اغلب لازم است؛ و \lr{offset} وقتی دقت زیاد در دامنه دینامیکی وسیع خواسته شود اهمیت پیدا می‌کند. اگر پس از برازش خطی، خطای کم‌بار همچنان سیستماتیک باشد، باید رجیسترهای \lr{xIRMSOS/xVRMSOS/xWATTOS} تراشه ارزیابی شوند.

\subsection{انتخاب نقاط و برداشت داده}
در هر نقطه، پس از زمان نشست، دست‌کم ۲۰ تا ۶۰ نمونه خام ثبت و میانگین شوند. ابزار مرجع باید هم‌زمان خوانده شود. صفر کامل برای کالیبراسیون \lr{RMS} مناسب نیست؛ راهنمای تراشه نیز \lr{offset} را با یک سیگنال کم ولی غیرصفر اندازه می‌گیرد.

\begin{longtable}{p{0.25\textwidth}p{0.62\textwidth}}
\toprule
کمیت & نقاط پیشنهادی \\
\midrule
\endhead
ولتاژ \lr{RMS} & پایین، نامی و بالای بازه واقعی؛ برای شبکه ۲۳۰ ولت، در صورت پشتیبانی منبع مثلاً ۱۸۰، ۲۳۰ و ۲۵۰ ولت. \\
جریان \lr{RMS} & کمترین جریان موردنیازِ معتبر، یک نقطه میانی و ۸۰ تا ۱۰۰ درصد جریان نامی طراحی؛ از حد مجاز \lr{CT/burden/ADE} عبور نشود. \\
توان اکتیو & بارهای مرجع کم، میانی و زیاد با ولتاژ پایدار و \lr{PF} معلوم؛ برای \lr{gain} ساده، بار مقاومتی نزدیک \lr{PF=1}. \\
توان ظاهری & هم‌زمان از ولتاژ و جریان مرجع یا پاورآنالایزر معتبر ثبت شود؛ صرف حاصل‌ضرب مقادیر نامطمئن استفاده نشود. \\
\bottomrule
\end{longtable}

برای بررسی خطای فاز \lr{CT} یک آزمون جدا با بار پایدار نزدیک \lr{PF=0.5} لازم است. نقاط برازش بهتر است کل بازه کاری را بپوشانند، ولی یک یا دو نقطه اضافی باید فقط برای اعتبارسنجی نگه داشته شوند و وارد برازش نشوند.

\subsection{کد کالیبراسیون ولتاژ، جریان و توان}
در نمونه زیر نام‌های \lr{raw\_...} میانگین شمارهای ثبت‌شده از
\lr{g\_app\_electrical.raw.phase[0]}
هستند و اعداد ستون دوم مقدار هم‌زمان ابزار مرجع در واحد مهندسی هستند:

\Needspace{26\baselineskip}
\begin{LTR}
\begin{lstlisting}[language=C]
const ADE7880_CalibrationPoint voltage_points[] = {
    {raw_v_180, 180.0F}, {raw_v_230, 230.0F},
    {raw_v_250, 250.0F}
};
const ADE7880_CalibrationPoint current_points[] = {
    {raw_i_0_10, 0.10F}, {raw_i_1_00, 1.00F},
    {raw_i_5_00, 5.00F}
};
const ADE7880_CalibrationPoint power_points[] = {
    {raw_p_20, 20.0F}, {raw_p_200, 200.0F},
    {raw_p_1000, 1000.0F}
};
ADE7880_LinearCalibration voltage_result;
ADE7880_LinearCalibration current_result;
ADE7880_LinearCalibration power_result;
float voltage_max_error_v;
float current_max_error_a;
float power_max_error_w;

(void)App_CalibratePhaseMultiPoint(
    APP_PHASE_QUANTITY_VOLTAGE_RMS, ADE7880_PHASE_A,
    voltage_points, 3U, &voltage_result, &voltage_max_error_v);
(void)App_CalibratePhaseMultiPoint(
    APP_PHASE_QUANTITY_CURRENT_RMS, ADE7880_PHASE_A,
    current_points, 3U, &current_result, &current_max_error_a);
(void)App_CalibratePhaseMultiPoint(
    APP_PHASE_QUANTITY_ACTIVE_POWER, ADE7880_PHASE_A,
    power_points, 3U, &power_result, &power_max_error_w);
\end{lstlisting}
\end{LTR}

سه مقدار \lr{max\_error} بیشترین فاصله نقاط ورودی از خط برازش‌شده را به‌ترتیب برحسب ولت، آمپر و وات نشان می‌دهند. کوچک‌بودن آن‌ها کافی نیست؛ خروجی باید در نقاط مستقل نیز با مرجع مقایسه شود. برای \lr{CT} معکوس، مقدار مرجع همچنان خروجی نهایی مطلوب است و تابع کاربردی علامت ضریب توان اکتیو را با تنظیم \lr{current\_polarity} هماهنگ می‌کند.

نتیجه نصب‌شده تا \lr{reset} در \lr{RAM} است. پس از تأیید، اعضای \lr{scale/offset} ساختارهای نتیجه در \lr{Watch} یادداشت و در \lr{Flash} یا تنظیمات \lr{startup} ذخیره شوند. نمونه بارگذاری ثابت تأییدشده:

\begin{LTR}
\begin{lstlisting}[language=C]
static const ADE7880_LinearCalibration saved_power_a = {
    .scale = VERIFIED_POWER_A_SCALE,
    .offset = VERIFIED_POWER_A_OFFSET
};

(void)App_SetPhaseCalibration(
    APP_PHASE_QUANTITY_ACTIVE_POWER,
    ADE7880_PHASE_A, &saved_power_a);
\end{lstlisting}
\end{LTR}

جریان نول نیز مسیر چندنقطه‌ای همسان و مستقل دارد:

\begin{LTR}
\begin{lstlisting}[language=C]
const ADE7880_CalibrationPoint neutral_points[] = {
    {raw_n_0_10, 0.10F},
    {raw_n_1_00, 1.00F},
    {raw_n_5_00, 5.00F}
};
ADE7880_LinearCalibration neutral_result;

(void)App_CalibrateNeutralCurrentMultiPoint(
    neutral_points, 3U, &neutral_result, NULL);
\end{lstlisting}
\end{LTR}

توان نوشته‌شده روی بدنه لامپ فقط مقدار نامی است و جای وات‌متر یا منبع کالیبره را نمی‌گیرد. ترتیب حرفه‌ای پیشنهادی این است: جهت \lr{CT}، بررسی فاز، کالیبراسیون \lr{gain} چندنقطه‌ای، ارزیابی نقطه‌های مستقل، و فقط در صورت باقی‌ماندن خطای کم‌بار، اصلاح \lr{offset} خود تراشه.

\section{آزمون سخت‌افزاری مرحله‌به‌مرحله}
به علت اتصال ورودی‌های \lr{ADE7880} به برق شهر، این آزمون فقط با ایزولاسیون مناسب، منبع محدودشده و فرد آگاه به ایمنی برق انجام شود. اتصال مستقیم اسیلوسکوپ یا دیباگر زمین‌شده به بخش غیرایزوله می‌تواند خطرناک و مخرب باشد.

\subsection{آزمون بدون اعمال ولتاژ اندازه‌گیری}
\begin{enumerate}
  \item بارهای رله و ورودی‌های قدرت را ایمن یا جدا کنید.
  \item اتصال \lr{SWDIO=PA13}، \lr{SWCLK=PA14}، زمین مشترک و مرجع ۳.۳ ولت \lr{J-Link} را برقرار کنید.
  \item برنامه را بسازید، پروگرام و \lr{reset} کنید.
  \item اگر \lr{ADE7880} پاسخ دهد، \lr{LED}ها باید الگوی رقص تقریباً یک‌ثانیه‌ای داشته باشند.
  \item اگر ارتباط برقرار نشود، دو \lr{LED} باید سریع یکی‌درمیان شوند و \lr{communication\_errors} پس از تلاش‌های خواندن افزایش یابد.
  \item با \lr{logic analyzer} کنترل کنید \lr{CS} در هر تراکنش پایین، \lr{SCLK} در بی‌کاری بالا و نرخ آن تقریباً ۵۶۲.۵ کیلوهرتز باشد.
  \item فرمان خواندن باید با \lr{0x01} و فرمان نوشتن با \lr{0x00} آغاز شود و سپس نشانی ۱۶بیتی \lr{MSB-first} بیاید.
\end{enumerate}

\subsection{آزمون با منبع مرجع}
\begin{enumerate}
  \item ابتدا ولتاژ کم و ایمن مطابق طراحی ورودی اعمال کنید و تغییر رجیستر متناظر را ببینید.
  \item فقط پس از تأیید سیم‌بندی، منبع \lr{AC} مرجع و بار کنترل‌شده وصل شود.
  \item فازهای \lr{A}، \lr{B} و \lr{C} جداگانه تغییر داده شوند تا نگاشت هر سه کانال قطعی شود.
  \item علامت توان اکتیو با جهت توان کنترل شود؛ برای بار مصرفی باید مثبت باشد. اگر اصلاح فیزیکی ممکن نیست، \lr{App\_SetCurrentPolarity} یک‌بار تنظیم شود.
  \item ضریب توان یک بار مقاومتی باید به مثبت یک نزدیک باشد. برای بار غیرخطی، \lr{power\_factor\_abs} بزرگی و \lr{power\_factor} اطلاعات علامت پیش‌فاز/پس‌فاز را نشان می‌دهد.
  \item قطع موقت سیم \lr{SPI} یا تغذیه \lr{ADE} باید حالت خطا ایجاد کند و پس از رفع عیب، تلاش دوره‌ای باید ارتباط را بازیابی کند.
  \item آزمون پیوسته حداقل ۲۴ ساعت برای شمار خطا، پایداری \lr{snapshot} و الگوی \lr{LED} انجام شود؛ آزمون چندروزه و چرخه دما برای ادعای صنعتی لازم است.
\end{enumerate}

\section{راهنمای عیب‌یابی}
\begin{longtable}{p{0.36\textwidth}p{0.54\textwidth}}
\toprule
نشانه & بررسی پیشنهادی \\
\midrule
\endhead
\lr{LED} خطا از ابتدای روشن‌شدن & تغذیه و کلاک کریستال \lr{ADE}، اتصال \lr{PM0/PM1}، سطح \lr{RESET}، چهار خط \lr{SPI}، سطح \lr{CS} و مشترک‌بودن زمین بررسی شود. \\
تمام خواندن‌ها \lr{0xFFFFFFFF} & مسیر \lr{MISO} باز، \lr{CS} غیرفعال، تراشه بی‌تغذیه یا رابط هنوز روی \lr{I2C} است. \\
تمام خواندن‌ها صفر & کلاک و ورودی‌ها، \lr{RUN}، اتصال کوتاه \lr{MISO} یا \lr{reset} دائمی بررسی شود. \\
خطای متناوب & طول سیم، نویز، زمین، تغذیه، خازن‌های بای‌پس و شکل موج \lr{SCLK/CS} با \lr{logic analyzer} بررسی شود. \\
\lr{PF} درست ولی ولت/آمپر نادرست & ابتدا داده \lr{raw} را کنترل و سپس نسبت ورودی و ضرایب کالیبراسیون را اصلاح کنید. \\
توان مصرفی منفی است & جهت \lr{CT} را در حالت بی‌برق با علامت‌های سازنده و شماتیک تطبیق دهید یا جبران نرم‌افزاری \lr{REVERSED} را یک‌بار اعمال کنید. \\
توان مثبت ولی \lr{PF} منفی است & الزاماً خطا نیست؛ \lr{PF} علامت‌دار \lr{ADE7880} پیش‌فاز/پس‌فاز را مشخص می‌کند. بزرگی را از \lr{power\_factor\_abs} بخوانید و با پاورآنالایزر مرجع مقایسه کنید. \\
ناشناخته‌بودن نوع \lr{DMA} در \lr{Keil} & فعال‌بودن ماژول \lr{DMA} در تنظیمات \lr{HAL} و عضویت سورس آن در پروژه بررسی شود. \\
پس از \lr{CubeMX} کد کامپایل نمی‌شود & عضویت فایل‌های برنامه و کتابخانه، مسیرهای \lr{include} و تعریف \lr{ADE\_CS} بررسی شود. \\
\bottomrule
\end{longtable}

در همین خطا وجود فراخوانی \lr{App\_Init(\&hspi2)} در بلوک کاربر فایل اصلی نیز کنترل شود.

\section{نتیجه اعتبارسنجی نرم‌افزاری این مرحله}
ساخت تازه با
\lr{PlatformIO ststm32 20.0.0}
و
\lr{framework-stm32cubef1 1.8.7}
با گزینه‌های
\lr{-Wall -Wextra -Werror}
موفق شد. خروجی ۱۰۳۷۲ بایت \lr{Flash} و ۵۴۴ بایت \lr{RAM} مصرف کرد. ساخت دوم با \lr{Keil} و \lr{Arm Compiler 6.24} نیز با \lr{0 Error(s), 0 Warning(s)} پایان یافت و اندازه‌های \lr{Code=9848}، \lr{RO-data=580}، \lr{RW-data=12} و \lr{ZI-data=1532} بایت را گزارش کرد. فایل‌های اجرایی هر دو زنجیره‌ابزار تولید شدند. جدول پیش‌فرض \lr{A/B/C} و نول، بازنشانی ضرایب و کالیبراسیون چندنقطه‌ای نول در هر دو زنجیره‌ابزار ساخته شدند. در ممیزی نهایی، تغییر ضریب یا جهت ترانس جریان نیز اصلاح شد تا مقدار مهندسی و نمایش قدیمی همان لحظه پاک شوند و دیباگر تا نمونه بعدی داده کهنه نشان ندهد. کاربر پیش‌تر برقراری ارتباط، دیباگ \lr{Keil}، مشاهده ولتاژ و پاسخ جریان/توان یک بار لامپی روی برد واقعی را گزارش کرده است. دقت ضرایب جایگزین موقت فاز \lr{C} و نول و دقت نهایی همه کانال‌ها هنوز باید با ابزار مرجع و نقاط مستقل روی برد کنترل شود.

\section{چک‌لیست پذیرش این قابلیت}
این مرحله زمانی از نظر سخت‌افزاری پذیرفته می‌شود که همه موارد زیر ثبت شوند:

\begin{itemize}
  \item \lr{build} بدون \lr{warning} در \lr{PlatformIO} و \lr{Keil}؛
  \item تولید مجدد \lr{CubeMX} و باقی‌ماندن منطق \lr{Application/Library}؛
  \item مشاهده کلاک، \lr{CS} و قالب فرمان صحیح با \lr{logic analyzer}؛
  \item \lr{online=true} و افزایش پیوسته شمار نمونه؛
  \item عدم افزایش خطای ارتباط در آزمون طولانی؛
  \item صحت نگاشت فازهای \lr{A}، \lr{B} و \lr{C} و کانال نول؛
  \item کالیبراسیون هر کانال با مرجع معتبر و ثبت ضرایب؛
  \item بازیابی خودکار پس از قطع و وصل کنترل‌شده ارتباط یا تغذیه؛
  \item ثبت نتیجه، تاریخ، نسخه \lr{commit} و تجهیزات آزمون در پرونده پروژه.
\end{itemize}
